\RR{The reviewer's objection to the bin factor is well taken, and the quantities in
\eqref{eq:slope} do not require bins at all.} The onsets of partial activations are
a point process on the delay window with intensity $\nu(t)$. Writing
$\ln\nu(t)=c_0+c_1t+c_2t^{2}$, the log likelihood of an inhomogeneous Poisson process
observed on $[T_0,T_1]$ is
\begin{equation}
\ell(c)=\sum_i \ln\nu(t_i)-\int_{T_0}^{T_1}\nu(t)\,\dd t,
\label{eq:ppll}
\end{equation}
a function of the exact event times, from which $c_1$ and $c_2$ are read off directly.
Three remarks.

\NN{Three estimands, not one.} They are not interchangeable, and this note uses all
three. A \emph{linear contrast} on the raw rate estimates $\dd\nu/\dd t$, an additive
slope, and is what Section~\ref{sec:reanalysis} reports. A \emph{Poisson regression with
a log link}, binned or not, estimates $\dd\ln\nu/\dd t$, a multiplicative slope, and is
what \eqref{eq:slope} predicts. The \emph{point-process likelihood} \eqref{eq:ppll}
is the unbinned form of the second, not of the first.

The two scales can disagree, and they do. During the stimulus period the coherence
effect on the slope is significant on the raw scale ($p=.007$) and absent on the log
scale ($p=.25$), because the log slope divides by the level and the high-coherence cell
has a much higher rate, so the same additive rise is a smaller relative rise. That is
not a power problem: the log slopes are non-monotone in coherence at short duration
($0.550$, $0.782$, $0.703$) where the raw contrasts are monotone. Neither scale is
privileged during the stimulus, because the model predicts no functional form there.
During the delay it predicts a log-linear one, which is why the point-process analysis
is restricted to that window.

\emph{Bins are harmless if they are an offset rather than a factor.} A Poisson
regression of the count in a bin of width $w$ on the bin centre, with $\log(w\times
n_{\text{trials}})$ as offset, is the discretised form of the same estimating equation,
because
$\E[\text{count}]=w\nu(t)\bigl[1+\tfrac{w^{2}}{24}\bigl((\ln\nu)''+((\ln\nu)')^{2}\bigr)+O(w^{4})\bigr]$;
at $w=0.1$~s and the values estimated here that is a relative error of
$2.5\times10^{-4}$ on the rate and none on the slope. The bin width is therefore not an
arbitrary analytic choice in the sense the reviewer objects to. What was arbitrary was
entering the bin as a nineteen-level factor.

\emph{The process is not Poisson, so the likelihood is a quasi-likelihood.} Upcrossings
cluster: an excursion that dips briefly below the gate and returns produces two events a
few milliseconds apart, which is also why the detection procedure merges them. Counts
are consequently overdispersed and the Poisson information understates the variance.
The Poisson likelihood is therefore best read as a quasi-likelihood, or an estimating
equation for the first-order intensity, rather than as the likelihood of the process.
Standard errors must come from between-participant variation. A further caveat belongs
here: what the detector delivers is the intensity of clusters, which is the upcrossing
intensity times a retention factor slightly below one, and that factor depends on the
distance to the gate and so drifts across the window. The measured log-rate slope is thus
the sum of the quantity of interest and the derivative of $\log$ of that factor, which is
small but not identically zero.

\emph{Clustering also enters the intensity.} What the detector delivers is the intensity
of clusters, which is the upcrossing intensity times a factor slightly below one. That
factor depends on the distance to the gate and so drifts across the window, contributing
a small bias to the estimated slope.

\paragraph{Results.} \NN{Two estimators are computed, in R, and they agree.} The
two-stage version fits the intensity separately for each participant and cell and enters
the coefficients into the same repeated-measures analysis as any other score, so that the
second level supplies the standard errors and the clustering does not bias inference. The
one-stage version fits a single log-linear model to all participants with cluster-robust
standard errors.

\begin{center}
\begin{tabular}{lcc}
\toprule
& two-stage & one-stage, robust\\
\midrule
log-rate slope & $+0.929$, $F(1,21)=102.95$, $p<.001$ & $+0.841$, $z=11.23$, $p<.001$\\
slope $\times$ duration & $F(1,21)=4.78$, $p=.040$ & $-0.179$, $z=-3.17$, $p=.002$\\
slope $\times$ coherence & $F(1.38,29.04)=4.20$, $p=.038$ & $-0.263$, $z=-5.02$, $p<.001$\\
slope $\times$ both & $F(1.87,39.23)=0.39$, $p=.663$ & ---\\
\bottomrule
\end{tabular}
\end{center}

The slope is positive, shallower after a long stimulus and shallower at high coherence
--- both of which follow from \eqref{eq:slope}, since either raises the maintained level
and so reduces the distance to the gate. The coefficients differ between the two estimators by about a tenth, which is expected: the
two-stage estimator weights participants equally and the one-stage weights them by
precision. A difference in sign would not have been expected.

